% !TeX program = pdflatex
\documentclass[11pt,a4paper]{article}
\usepackage[T1]{fontenc}
\usepackage[utf8]{inputenc}
\usepackage{lmodern,amsmath,amssymb,amsthm,mathtools,esint}
\usepackage[margin=25mm,headheight=14pt]{geometry}
\usepackage{microtype,booktabs,array,enumitem,xurl}
\usepackage[unicode,hidelinks]{hyperref}
\usepackage{fancyhdr}
\pagestyle{fancy}\fancyhf{}\fancyhead[L]{Four-dimensional Einstein geometry}\fancyhead[R]{8 October 2026}\fancyfoot[C]{\thepage}
\setlength{\emergencystretch}{3em}
\setlist{nosep,leftmargin=*}
\allowdisplaybreaks[2]
\numberwithin{equation}{section}
\newtheorem{theorem}{Theorem}[section]
\newtheorem{proposition}[theorem]{Proposition}
\newtheorem{corollary}[theorem]{Corollary}
\newtheorem{lemma}[theorem]{Lemma}
\theoremstyle{definition}\newtheorem{hypothesis}[theorem]{Input}
\theoremstyle{remark}\newtheorem{remark}[theorem]{Remark}
\newcommand{\Ric}{\operatorname{Ric}}
\newcommand{\Rm}{\operatorname{Rm}}
\newcommand{\tr}{\operatorname{tr}}
\newcommand{\Met}{\operatorname{Met}}
\newcommand{\CK}{\operatorname{CK}}
\newcommand{\CP}{\mathbb{CP}}
\newcommand{\Sym}{\operatorname{Sym}}
\newcommand{\kerop}{\operatorname{ker}}
\newcommand{\imop}{\operatorname{im}}
\title{Four Questions in Einstein Geometry:\\[3pt]
Conformal-K\"ahler Separation, the Einstein Weyl Channel,\\
and the Complete Darboux Ricci Symbol}
\author{Mathematical continuation of the Qiushi 8 October handoff}
\date{8 October 2026}
\begin{document}
\maketitle
\begin{abstract}
We use the corrected statements in the two-session handoff of 8 October.
The background-sensitive strict Burns saving is an input of that handoff.
A scalar-energy argument extends the comparison barrier to all globally
conformally K\"ahler metrics, without fixing their complex structure, and
separates subthreshold minimizing sequences from that entire class.
A concrete Horikawa blowup removes Einstein nonexistence from the list of
unproved requirements for one possible minimizer example.
An exact inversion of the recorded classical Weyl bound identifies the
Weyl deficit forced on an Einstein competitor; it is a reformulation of the
classical envelope, not an improved unrestricted lower bound.
Finally, we calculate the full image, kernel, gauge coercivity, and a right
inverse of the Darboux-compatible Ricci principal symbol.  The kernel is
exactly the Hamiltonian gauge direction.  These statements provide concrete
inputs for the four questions, but neither smooth global attainment nor the
negative Goldberg conclusion follows from them.
\end{abstract}

\section{The four questions and the source inputs}

The present questions are those of \cite[Section 1]{Synthesis}:
Einstein existence on $M=X\#k\overline{\CP}^{\,2}$ with $X$ minimal of
general type and $k\ge1$; the unrestricted infimum of $P$; a smooth closed
non-Bach-flat global minimizer of $P$; and compact negative Einstein
almost-K\"ahler integrability in dimension four.  The older O10/O11 labels
are not substituted for these four specifications.

On a closed oriented four-manifold set
\begin{equation}
 P(g)=\int_M|\Ric_g^\circ|^2\,dV_g,\qquad
 S(g)=\int_M s_g^2\,dV_g,\qquad
 W_\pm(g)=\int_M|W_g^\pm|^2\,dV_g,
 \quad T_\pm=2\chi\pm3\tau.
\end{equation}
Weyl norms are Hilbert--Schmidt norms on two-forms, while Ricci is a
symmetric-tensor norm.  The identities used throughout are
\begin{equation}\label{eq:CW}
 P=\frac S{12}+4W_\pm-8\pi^2T_\pm,
 \qquad \int_M|\Rm_{\rm op}|^2=8\pi^2\chi+P.
\end{equation}
Write $p_*(M)=\inf_{g\in\Met(M)}P(g)$; the normalized relative Einstein
functional is not denoted by $p_*$.

For a minimal general-type model $X$ let $A=c_1^2(X)>0$.
The classical scalar budget on its blowup is
\begin{equation}\label{eq:scalar}
 S(g)\ge32\pi^2A\quad\text{for every smooth metric }g.
\end{equation}
This follows, for example, from the classical scalar-curvature invariant
of a general-type surface \cite[Introduction]{LeBrun}.
Applications below may instead take \eqref{eq:scalar} as their explicit
input.  The negative Yamabe property, when used, is stated separately.

The new overnight manuscript proves, in its stated geometric range,
\begin{equation}\label{eq:source-saving}
 p_*(X\#k\overline{\CP}^{\,2})<8\pi^2 k.
\end{equation}
Its mechanism is the optimized Burns response
\begin{equation}
 P(g_\varepsilon)=P(g)+8\pi^2
 -8\pi^2\varepsilon\mathbf R^+_p(\zeta,\zeta)
 +4\pi^2\Phi\varepsilon^2\log(1/\varepsilon)+O(\varepsilon^2),
\end{equation}
with $\Phi=-|T|^2$ on an Einstein kernel direction.  In the KE case the
coefficient is $-s^2/32$ before multiplication by $4\pi^2$.
We inherit \eqref{eq:source-saving} and its complete-support, background,
matching, and resolution hypotheses from \cite[Section 3]{Synthesis},
\cite[A01.52 and corrected A02.59D]{Companion}.  None of the arguments below
reproves the logarithmic construction.

\section{An energy barrier for the entire conformally K\"ahler class}

Let $\CK_+(M)$ be the set of smooth metrics globally conformal to a smooth
K\"ahler metric inducing the chosen orientation of $M$.  The complex
structure and the positive conformal factor are both allowed to vary.
Define $\CK_-(M)$ using the opposite orientation.  This definition does not
include merely locally conformally K\"ahler metrics or singular conformal
factors.

\begin{theorem}\label{thm:CK}
Assume \eqref{eq:scalar}.  If $g=u^2h$, with $u>0$ smooth and $h$ a smooth
K\"ahler metric of orientation $\epsilon\in\{+,-\}$, then
\begin{equation}\label{eq:CKexact}
 \boxed{P(g)-8\pi^2(A-T_\epsilon)
 =\frac{S(g)-32\pi^2A}{12}
  +\frac{S(h)-32\pi^2A}{6}\ge0.}
\end{equation}
In particular, on $M=X\#k\overline{\CP}^{\,2}$ with its complex orientation,
\begin{equation}\label{eq:CKbarrier}
 \inf_{g\in\CK_+(M)}P(g)\ge8\pi^2 k.
\end{equation}
If additionally $\tau(M)\ge0$, then
\begin{equation}\label{eq:bothCK}
 \inf_{g\in\CK_+(M)\cup\CK_-(M)}P(g)\ge8\pi^2k.
\end{equation}
\end{theorem}

\begin{proof}
For a K\"ahler surface with its complex orientation,
$|W^+_h|^2=s_h^2/24$.  This pointwise identity is classical
\cite[Introduction]{LeBrun}.  The integrated chiral Weyl energy is
conformally invariant in dimension four.  Consequently
\[
 W_\epsilon(g)=W_\epsilon(h)=S(h)/24.
\]
Substitution in \eqref{eq:CW} gives
\[
 P(g)=S(g)/12+S(h)/6-8\pi^2T_\epsilon.
\]
Subtract $8\pi^2(A-T_\epsilon)$ and use \eqref{eq:scalar} for both metrics.
This proves \eqref{eq:CKexact} without any assumption on how their complex
structure relates to the original one on $X$.
For a blowup, $T_+=A-k$.  If $\tau(M)\ge0$, then $T_-\le T_+$, so the
opposite-orientation lower bound is at least as large.
\end{proof}

\begin{corollary}[Separation of genuine minimizing sequences]\label{cor:separation}
Assume the input \eqref{eq:source-saving} and let $g_{\rm test}$ be a fixed
smooth competitor satisfying
\[
 P(g_{\rm test})\le8\pi^2k-\delta_0,\qquad \delta_0>0.
\]
Every minimizing sequence $P(g_j)\to p_*(M)$ eventually satisfies
\begin{equation}
 P(g_j)\le8\pi^2k-\delta_0/2
\end{equation}
and therefore lies outside $\CK_+(M)$.  If $\tau(M)\ge0$, it eventually
lies outside both conformally K\"ahler classes.  The same exclusion applies
to a smooth global minimizer whenever one exists.
\end{corollary}

\begin{proof}
The test metric gives $p_*(M)\le8\pi^2k-\delta_0$.  Convergence of the
energies gives the displayed bound for large $j$, while Theorem
\ref{thm:CK} excludes membership in the indicated classes.
\end{proof}

The separation is visible in a conformal invariant.  For any smooth $g$
with $P(g)\le8\pi^2k-\delta$, $\delta>0$, equation \eqref{eq:CW} and the
scalar budget imply
\begin{equation}\label{eq:Weylseparation}
 \boxed{W_+(g)\le\frac{4\pi^2A}{3}-\frac\delta4,
 \qquad W_+(h)\ge\frac{4\pi^2A}{3}\quad(h\in\CK_+(M)).}
\end{equation}
Thus changing a positive conformal factor cannot turn the low-cost metric
into a K\"ahler one.  This is a separation in Weyl energy, not a claimed
distance between metrics in a prescribed metric-space topology.

For a complex-hyperbolic $X$, $\tau(X)=A/9$, so the two-orientation statement
applies when $k\le A/9$.  The topology formulas are
\[
 T_+(M)=A-k,\qquad T_-(M)=A/3+5k,\qquad \tau(M)=A/9-k.
\]
The $+5k$ sign in $T_-$ is essential.

\paragraph{What this contributes to Q2 and Q3.}
The handoff's strict upper bound and Theorem \ref{thm:CK} give an energy
separation from the entire indicated conformally K\"ahler locus, not merely
from a fixed K\"ahler cone.  The theorem is a short consequence of classical
identities and the scalar budget; it is not represented as a new analytic
coercivity theorem or as a proof of smooth attainment.  No statement about
arbitrary Hermitian or locally conformally K\"ahler metrics is inferred.

\section{Exact Einstein candidates must overshoot the nominal Weyl threshold}

In the source's negative Yamabe representative, put
\begin{equation}
 \kappa=\frac A2-\frac{3W_+}{8\pi^2},\qquad
 \upsilon=\frac{S}{32\pi^2}-A,\qquad p=\frac P{8\pi^2}.
\end{equation}
These are exactly the Weyl-channel variables of
\cite[Section 7]{Synthesis}; $\kappa$ is not the O10 parameter defect.
We use the two recorded relations
\begin{align}
 p&=k-\frac43\kappa+\frac13\upsilon,\label{eq:channel}\\
 2\sqrt{A+\upsilon}+\sqrt{A-2\kappa}&\ge3\sqrt A.
 \label{eq:classical}
\end{align}
Dropping the nonnegative extra projection budget only weakens the second
inequality.  Its source is the classical combined scalar--Weyl estimate
\cite{LeBrun}, not a new lower bound.

\begin{proposition}[Inverting the classical feasible region]\label{prop:overshoot}
Let $\ell=k-p$ and suppose $0<\ell<A/3$.  Every competitor satisfying
\eqref{eq:channel}--\eqref{eq:classical} obeys
\begin{equation}\label{eq:overshoot}
 \boxed{\kappa\ge r_-(\ell)
 :=\frac{2A+6\ell-2\sqrt{A^2-3A\ell}}9
 =\ell+\frac{\ell^2}{A(1+\sqrt{1-3\ell/A})^2}.}
\end{equation}
Consequently, an Einstein metric in the range $0<k<A/3$ must have
\begin{equation}\label{eq:Einsteinkappa}
 \boxed{\kappa\ge k+
 \frac{k^2}{A(1+\sqrt{1-3k/A})^2}>k.}
\end{equation}
\end{proposition}

\begin{proof}
If $\kappa\le0$, then $\upsilon\ge0$ and \eqref{eq:channel} give $p\ge k$,
contrary to $\ell>0$.  For $0\le\kappa\le A/2$, solving
\eqref{eq:classical} for $\upsilon$ gives the source bound
\[
 p\ge L_A(k,\kappa)
 =k-\kappa+\frac12\{A-\kappa-\sqrt{A^2-2A\kappa}\}.
\]
This function is strictly decreasing up to $\kappa=4A/9$ and then
increasing.  Its smaller intersection with $p=k-\ell$ solves
\[
 \sqrt{A^2-2A\kappa}=A-3\kappa+2\ell,
\]
whose smaller root is the first expression for $r_-$.  For
$0<\ell<A/3$, this root lies in $(0,4A/9)$ and the unsquared right side is
positive there.  Thus $\kappa<r_-(\ell)$ is impossible.  Any $\kappa>4A/9$
already satisfies the desired lower bound.  Writing
$t=\sqrt{1-3\ell/A}$ gives
\[
 r_-(\ell)-\ell=\frac A9(1-t)^2
 =\frac{\ell^2}{A(1+t)^2},
\]
which proves the second expression.  For Einstein $P=0$, so $\ell=k$.
\end{proof}

This calculation fixes a precise target for an \emph{Einstein-specific}
upper estimate: it would be enough to prove $\kappa<r_-(k)$ on those exact
Einstein candidates.  Such an upper estimate has not been obtained here.
The stronger requirement $\kappa\le k$ for \emph{all} metrics would imply
it but is not necessary for Q1.  This distinction avoids making a full Q2
lower-bound theorem a prerequisite for every approach to Q1.

The restriction is still only classical feasibility.  For example, for
$0<k/A\le1/4$ the formal data
\begin{equation}
 P=0,\qquad \kappa=2k,\qquad \upsilon=5k
\end{equation}
satisfy both \eqref{eq:channel} and \eqref{eq:classical}.  To check the latter,
write $h=k/A$; it is
$2\sqrt{1+5h}+\sqrt{1-4h}\ge3$, which, after one legitimate squaring,
reduces to $\sqrt{1-4h}\ge1-4h$.
They are not a metric or an actual-field counterexample.  They show exactly
why inversion of the old budget cannot alone answer Q1.

\section{The attainment requirement in Q3}

Put $\mathcal W(g)=W_+(g)+W_-(g)$ and let $\mathfrak B$ denote its
$L^2$ gradient, with respect to variations of the covariant metric.  It is
a fixed nonzero multiple of the Bach tensor.  This definition avoids mixing
normalizations of the Bach tensor across references.

From \eqref{eq:CW},
\begin{equation}\label{eq:PtotalW}
 P=2\mathcal W+S/12-16\pi^2\chi.
\end{equation}
For a conformal variation $g_t=e^{2tf}g$ in dimension four,
\begin{equation}\label{eq:conformalP}
 \frac{d}{dt}\bigg|_{t=0}P(g_t)=-\int_M f\,\Delta s\,dV,
\end{equation}
where $\Delta=\nabla^i\nabla_i$.  Thus a smooth closed critical metric has
constant scalar curvature.  At such a metric,
$\nabla S=-2s\Ric^\circ$; hence criticality of \eqref{eq:PtotalW} gives
\begin{equation}\label{eq:Bach}
 \boxed{\mathfrak B=\frac{s}{12}\Ric^\circ.}
\end{equation}
These identities follow as well from the full quadratic-curvature
variation formulas in \cite[Section 2]{GV}.

On a manifold with negative Yamabe invariant a constant-scalar metric has
$s<0$.  Equation \eqref{eq:Bach} then yields
\begin{equation}
 \text{smooth $P$-critical and Bach-flat}\quad\Longleftrightarrow\quad
 \text{Einstein}.
\end{equation}
Therefore a valid route to a Q3 example on one chosen $M$ is
\[
 \text{smooth attainment on $M$}
 +\text{Einstein nonexistence on $M$}
 +\mathcal Y(M)<0.
\]
Theorem \ref{thm:CK} excludes conformally K\"ahler competitors below the
standard barrier, but it proves neither of the first two requirements.
Anderson's selected-minimizer framework \cite{Anderson} does not turn
orbifold or partly collapsed convergence into smooth attainment on the
original differentiable manifold without the further connections specified
in \cite[Section 4]{Synthesis}.


\subsection{A concrete candidate that does not depend on unresolved small-blowup nonexistence}

For an existence question such as Q3, the background need not be a large
ball quotient whose one-point blowup lies in the unresolved Q1 range.
The following classical surface gives a concrete alternative.
Let
\begin{equation}\label{eq:Horikawa}
 X=\{w^2+z^5+x^{10}+y^{10}=0\}
 \subset\mathbb P(1,1,2,5),\qquad M=X\#\overline{\CP}^{\,2}.
\end{equation}
The coordinates have weights $(1,1,2,5)$.  The affine cone is smooth away
from the origin, since the four partial derivatives vanish simultaneously
only there.  The hypersurface avoids the two singular points of the weighted
projective space, and hence is a smooth surface.  Weighted adjunction gives
$K_X=\mathcal O_X(1)$, an ample line bundle, with
\[
 K_X^2=\frac{10}{1\cdot1\cdot2\cdot5}=1,\qquad p_g(X)=2,\qquad q(X)=0.
\]
These are the classical $(1,2)$ Horikawa surfaces; their weighted
hypersurface description and invariants are recalled in
\cite[Introduction]{CP}.  In particular $X$ has a negative KE metric, and
its unperturbed monopole supply is in the $b^+>1$ range.
Noether's formula and the signature formula give
\begin{equation}
 \chi(X)=35,\quad\tau(X)=-23,\qquad
 \chi(M)=36,\quad\tau(M)=-24,\quad T_+(M)=0.
\end{equation}
The all-metric scalar and combined scalar--Weyl budgets for the blowup
use $A=1$.  The classical parameter inequality at $\beta=8$ is
\[
 D_8=\frac{27}{2}P+108\pi^2T_+-72\pi^2p_{\rm SW}\ge0,
 \qquad p_{\rm SW}\ge1.
\]
Thus the source's general-KE strict upper bound gives the explicit interval
\begin{equation}\label{eq:HorikawaWindow}
 \boxed{\frac{16\pi^2}{3}\le p_*(M)<8\pi^2.}
\end{equation}
The lower bound is classical.  The strict upper bound is the overnight
construction \eqref{eq:source-saving}, applied to this fixed smooth
negative KE background.  This combination does not determine $p_*(M)$.

Since $p_*(M)>0$, this $M$ carries no Einstein metric.  It also has negative
Yamabe invariant by the standard general-type scalar theory.  Consequently
\begin{equation}\label{eq:concrete-Q3}
 \boxed{\text{If $P$ attains its infimum smoothly on this $M$, then the
 minimizer is non-Bach-flat.}}
\end{equation}
Theorem \ref{thm:CK} additionally excludes globally conformally K\"ahler
minimizers inducing the displayed orientation.  No opposite-orientation
conclusion is asserted here, since $\tau(M)<0$.

This choice removes the unresolved Q1 Einstein alternative from this one
Q3 attempt.  What remains is a genuine attainment problem.  The strict
upper bound excludes a separate standard Burns loss of energy $8\pi^2$
under the stated additive-energy hypotheses, but does not exclude cheaper
ALE channels, Ricci-flat bubbles, or partial collapse.  Therefore
\eqref{eq:concrete-Q3} is not an example of an attained minimum.  Its value
is to isolate attainment on a concrete smooth manifold using already
available nonexistence and upper-bound information.

\section{The complete Darboux Ricci principal symbol}

This section supplies a concrete replacement for an ill-conditioned local
jet block.  It works at a single point and nonzero real cotangent vector;
it does not assert formal integrability to all orders.

Let $(V,g,J)$ be a Hermitian real four-space.  Identify a symmetric tensor
$H$ with its $g$-symmetric endomorphism and set
\[
 \mathcal E_J=\{H:H J+JH=0\},\qquad \dim_\mathbb R\mathcal E_J=6.
\]
For a real covector $\xi$, identified with its dual vector, define the Ricci
symbol in the source's convention by
\begin{equation}\label{eq:Ric-symbol}
 S_\xi(H)=\frac12\{
 \xi\otimes H\xi+H\xi\otimes\xi-|\xi|^2H\}.
\end{equation}
A simultaneous change of sign does not affect the statements about
kernel, image, or norms.

\begin{theorem}\label{thm:symbol}
For every $\xi\ne0$,
\begin{align}
 \kerop S_\xi&=\mathbb R(\xi\otimes J\xi+J\xi\otimes\xi),\label{eq:kernel}\\
 \imop S_\xi&=\left\{K\in\Sym^2V^*:
 K\xi-\tfrac12(\tr K)\xi=0,\quad
 K(J\xi,J\xi)-\tfrac12|\xi|^2\tr K=0\right\}.
 \label{eq:image}
\end{align}
In particular the image has dimension five: the four contracted-Bianchi
conditions and the additional scalar Darboux condition are the complete
linear conditions at this fixed real covector.
Moreover
\begin{equation}\label{eq:coercivity}
 \boxed{|S_\xi(H)|^2+\frac14|H(\xi,J\xi)|^2
 \ge\frac18|\xi|^4|H|^2.}
\end{equation}
The constant $1/8$ is sharp.
\end{theorem}

\begin{proof}
By unitary covariance and homogeneity, it is enough to take $\xi=e_1$,
$Je_1=e_2$, $Je_3=e_4$.  Every $H\in\mathcal E_J$ has the form
\begin{equation}\label{eq:Hmatrix}
 H=\begin{pmatrix}
 a&b&c&d\\ b&-a&d&-c\\ c&d&e&f\\ d&-c&f&-e
 \end{pmatrix}.
\end{equation}
Direct substitution gives
\begin{equation}\label{eq:Smatrix}
 S_{e_1}(H)=\frac12\begin{pmatrix}
 a&0&0&0\\0&a&-d&c\\0&-d&-e&-f\\0&c&-f&e
 \end{pmatrix}.
\end{equation}
The only absent variable is $b$, which proves \eqref{eq:kernel} and the
rank-five statement.  Conversely, the equations on the right of
\eqref{eq:image} force
$K_{12}=K_{13}=K_{14}=0$, $K_{11}=K_{22}$, $K_{44}=-K_{33}$.
These leave precisely the five entries in \eqref{eq:Smatrix}, so there is
no further image condition at this covector.
Finally,
\begin{align*}
 |S_{e_1}(H)|^2&=\tfrac12(a^2+c^2+d^2+e^2+f^2),\\
 |H|^2&=2(a^2+b^2+e^2+f^2)+4(c^2+d^2),\\
 H(e_1,Je_1)&=b.
\end{align*}
Consequently the difference between the left side and the right side of
\eqref{eq:coercivity} is exactly
$\tfrac14(a^2+e^2+f^2)\ge0$.  Equality holds, for example, when only $c$
is nonzero.  This proves sharpness and the theorem.
\end{proof}

\subsection{An explicit image projection and right inverse}
For a general symmetric matrix $K$ at $\xi=e_1$, its orthogonal projection
onto \eqref{eq:image} is
\begin{equation}
 \mathbb P K=
 \begin{pmatrix}
 t&0&0&0\\0&t&K_{23}&K_{24}\\
 0&K_{23}&u&K_{34}\\0&K_{24}&K_{34}&-u
 \end{pmatrix},\qquad
 t=\frac{K_{11}+K_{22}}2,\quad u=\frac{K_{33}-K_{44}}2.
\end{equation}
The gauge-fixed right inverse $\mathbb R$ on the image is obtained in
\eqref{eq:Hmatrix} by
\begin{equation}
 a=2K_{11},\quad b=0,\quad c=2K_{24},\quad
 d=-2K_{23},\quad e=-2K_{33},\quad f=-2K_{34}.
\end{equation}
It satisfies $S_{e_1}\mathbb R K=K$ on the image and
$|\mathbb R K|\le\sqrt8|K|$.
For unit $\xi$, the squared singular values of $S_\xi$ are
\begin{equation}
 0,\quad \tfrac18,\quad\tfrac18,\quad
 \tfrac14,\quad\tfrac14,\quad\tfrac14.
\end{equation}
These formulas permit explicit treatment of the image and gauge rather
than repeatedly inferring them from a nearly singular numerical Jacobian.

\subsection{Hamiltonian meaning and the nonlinear metric parametrization}
For a fixed Darboux symplectic form, Hamiltonian vector fields have symbol
$X_f=J\xi f$; the metric variation $\mathcal L_{X_f}g$ has principal part
$\xi\otimes J\xi+J\xi\otimes\xi$, up to the Fourier sign.  Thus the kernel
in \eqref{eq:kernel} is genuine gauge freedom, not evidence of a strictly
almost-K\"ahler Einstein deformation.  A scalar gauge with principal symbol
$H(\xi,J\xi)$ removes it; its action on the Hamiltonian potential has a
nonzero multiple of $|\xi|^4$ as symbol.

There is also a direct way to remove the algebraic metric constraints.
In a fixed Darboux frame choose $J_0^2=-I$ and use
\begin{equation}
 G=\exp H,\qquad H^T=H,\qquad HJ_0+J_0H=0.
\end{equation}
Then $G$ is positive definite, $\det G=1$, and
\begin{equation}
 GJ_0G=J_0.
\end{equation}
The compatible metric constraint and closure of the fixed symplectic form
are exact; they need not be reintroduced as soft numerical residuals.
This supplies six unconstrained coefficient functions before Hamiltonian
gauge fixing.  It is a local parametrization, not an imposed global
trivialization of the tangent bundle.

\subsection{What the symbol calculation does not imply}
The extra Darboux compatibility relation and the correction of the degree-15
count are already in \cite[A01.55]{Companion}.  Theorem \ref{thm:symbol}
adds the full fixed-covector kernel/image calculation, explicit right
inverse and norm control.  It does not show that a nonlinear lower-jet
source lies in the prolonged image, that the system is involutive, or that
a formal solution converges.  Curvature sign is absent from this principal
symbol; an obstruction distinguishing negative Einstein data from KE and
Ricci-flat controls must enter lower-order resolved compatibility or global
geometry.  The classical formal-integrability perspective is discussed in
\cite[Section 4]{ADsurvey}.


\subsection{A reproducible highest-order block calculation}

The accompanying script constructs the integer matrix of the homogeneous
highest-order operator $2D\Ric_\delta$ from degree-$N$ compatible tensor
polynomials to degree-$(N-2)$ symmetric tensor polynomials.  A run in this
turn gave the following ranks over both primes $65521$ and $65519$:
\begin{center}
\begin{tabular}{rrrrr}
\toprule
$N$ & Rows & Columns & Rank modulo each prime & Kernel minus Hamiltonian count\\
\midrule
4 &100&210&83&43\\
5 &200&336&156&60\\
6 &350&504&260&79\\
7 &560&720&400&100\\
\bottomrule
\end{tabular}
\end{center}
These agree with the corrected count $N^2+8N-5$ already recorded in
\cite[A01.55]{Companion}.  The degree-seven integer matrix occupies
3,225,600 bytes in this implementation.  Agreement modulo two primes is a
check of these finite matrices; it is not alone a proof of the rational
rank in every degree or of nonlinear prolongation compatibility.

The computation is deliberately different from solving the complete
nonlinear finite-jet equations: the known linear image, Hamiltonian gauge,
and compatibility conditions are separated first.  At each degree a new
nonlinear forcing term must still lie in the prolonged image.  Testing that
membership, and calculating any resulting lower-order condition, cannot be
replaced by the fixed-covector inverse above.

\section{Consequences for the four research tasks}

The following dependencies keep the four questions distinct.
\begin{enumerate}
\item A positive unrestricted lower bound for $P$ on one $M$ excludes its
Einstein metrics.  Q1 can also be approached directly at $P=0$, without
first computing $p_*(M)$.  Proposition \ref{prop:overshoot} gives the exact
classical Weyl threshold that such an argument must improve.
\item The strict upper bound in the handoff, together with Theorem
\ref{thm:CK}, places low-energy competition outside an entire conformally
K\"ahler class.  A matching lower bound remains necessary to determine
$p_*$; a conformal modification cannot remove this competition.
\item Q3 requires smooth attainment and exclusion of the Einstein
alternative on the same smooth manifold.  For the concrete manifold
\eqref{eq:Horikawa}, the Einstein alternative is already excluded and
\eqref{eq:HorikawaWindow} gives an explicit energy window.  Smooth attainment
remains unproved; the Bach identity is an implication, not an existence theorem.
\item Q4 still needs a new resolved global or lower-order compatibility
constraint, or a complete closed counterexample.  The complete symbol
calculation changes the numerical formulation and prevents counting gauge
and known identities as mathematical obstructions.
\end{enumerate}

No result here declares Q1, Q2, Q3, or Q4 fully solved.  The additions are the
complete conformally K\"ahler separation argument, the explicit inversion of
the source's classical envelope, and the gauge-resolved Ricci-symbol
calculation.  The general-background negative logarithmic construction
remains the system's result, and the scalar, curvature-variation and symbol
background remains classical.  These provenance distinctions are part of
the mathematical statement, not separate research objectives.

\begin{thebibliography}{9}
\bibitem{Synthesis}
Qiushi Sessions S1003A01 and S1003A02,
\emph{Einstein Gaps on Four-Manifolds: Four Public Problems},
8 October 2026, supplied file \texttt{01\_Research\_Synthesis(1).tex}.
The sections on the four tasks, the insertion theorem, minimizers,
Weyl coordinates, and Goldberg compatibility specify the four tasks, the corrected Burns theorem,
attainment conditions, Weyl-channel bounds, and Goldberg corrections.
\bibitem{Companion}
Same sessions, \emph{Mathematical Source Companion}, supplied file
\texttt{02\_Mathematical\_Source\_Companion(1).tex}.
Source keys A01.52, A01.55, A02.58, A02.59D, A02.62--64.
\bibitem{LeBrun}
C. LeBrun, \emph{Ricci curvature, minimal volumes, and Seiberg--Witten theory},
Invent. Math. 145 (2001), 279--316.
\url{https://arxiv.org/abs/math/0003068}.
\bibitem{CP}
S. Coughlan and R. Pignatelli,
\emph{Simple fibrations in $(1,2)$-surfaces},
Forum of Mathematics, Sigma (2023).
\url{https://doi.org/10.1017/fms.2023.41}.
The introduction recalls the classical Horikawa surface description;
no new surface-existence theorem is claimed here.
\bibitem{GV}
M. J. Gursky and J. A. Viaclovsky,
\emph{Rigidity and stability of Einstein metrics for quadratic curvature
functionals}, arXiv:1105.4648v2 (2013).
\url{https://arxiv.org/abs/1105.4648}.
\bibitem{Anderson}
M. T. Anderson, \emph{Canonical metrics on 3-manifolds and 4-manifolds},
Asian J. Math. 10 (2006), 127--163.
\url{https://arxiv.org/abs/math/0511034}.
\bibitem{ADsurvey}
V. Apostolov and T. Dr\u aghici,
\emph{The curvature and the integrability of almost-K\"ahler manifolds:
a survey}, 2003.
\url{https://arxiv.org/abs/math/0302152}.
\bibitem{AAD}
V. Apostolov, J. Armstrong and T. Dr\u aghici,
\emph{Local models and integrability of certain almost K\"ahler 4-manifolds},
Math. Ann. 323 (2002), 633--666.
\url{https://arxiv.org/abs/math/0109098}.
\end{thebibliography}
\end{document}
